%%%%%%%%%%%%%%%%%%%%%%%%%%%%%%%%%%%%%%%%%%%%%%%%%%%%%%%%%%%%%%%%%%%%%%%%%%%
%
% HST Cycle 34 Phase I Proposal --- ABSTRACT
%
% Title: Catching IGM Dust in the Act:
%        A Per-Sightline UV Spectroscopic Test with HST/COS and LSST SNe Ia
%
% This file contains the APT-submitted abstract (plain-text friendly).
% It is written as a standalone compilable LaTeX snippet so the PI can
% (a) build it on its own to preview, and (b) paste the rendered text
% directly into the APT abstract field.
%
%%%%%%%%%%%%%%%%%%%%%%%%%%%%%%%%%%%%%%%%%%%%%%%%%%%%%%%%%%%%%%%%%%%%%%%%%%%

\documentclass[11pt]{article}
\usepackage[letterpaper,margin=1in]{geometry}
\usepackage{amsmath}
\emergencystretch=1em

% Light-weight macros so the file compiles standalone without the
% HST style file (which defines \ion elsewhere). Convert the numeric
% second argument to an uppercase Roman numeral so "\ion{Mg}{2}"
% renders and extracts as "Mg II".
\makeatletter
\newcommand{\@ROMANNUM}[1]{%
  \ifcase#1\or I\or II\or III\or IV\or V\or VI\or VII\or VIII\or IX\or X\fi}
\providecommand{\ion}[2]{#1\,\@ROMANNUM{#2}}
\makeatother

\begin{document}

\begin{center}
{\large\bfseries Catching IGM Dust in the Act:\\
A Per-Sightline UV Spectroscopic Test with HST/COS and LSST SNe~Ia}
\end{center}

\noindent
Type~Ia supernovae (SNe~Ia) are the cornerstone of modern cosmology, and
in the Stage~IV era systematic uncertainties now dominate their distance
budget. The Union3 analysis identifies intergalactic dust extinction as
the second-largest systematic in SNe~Ia cosmology --- yet the existence
of this dust has only ever been inferred \emph{statistically}, through
the cross-correlation of quasar colors with foreground galaxies
(M\'enard et~al.\ 2010). In the sixteen years since, no study has
obtained direct spectroscopic evidence that reddening along any
individual SNe~Ia sightline is caused by intervening \ion{Mg}{2}
absorption.

\begin{sloppypar}
We propose a 34-orbit non-disruptive Target of Opportunity program to
obtain HST/COS~G185M UV spectroscopy of 5--7 anomalously reddened
SNe~Ia in the COSMOS Deep Drilling Field, triggered within three weeks
of LSST/Rubin discovery. A five-stage pipeline --- LSST alert; COSMOS2020
photo-$z=0.30$--0.45; dDLR$>3$ (HST-ACS morphology) for host-dust
rejection; foreground-galaxy impact parameter $\rho<200$~kpc; and LSST
SALT3 color excess $\Delta E(B{-}V)\geq0.03$~mag --- yields 6--10
COS-ready targets per year. At the redshifts of the pre-identified
foreground galaxies we search for
\ion{Mg}{2}~$\lambda\lambda$2796,~2803 absorption; for sightlines
placing Ly$\alpha$ in COS~G130M, we conduct a bonus DLA search.
\end{sloppypar}

A single detection correlated with the photometric reddening would
constitute the \emph{first direct, per-sightline spectroscopic evidence
in history} that IGM dust biases the SNe~Ia Hubble diagram ---
converting a 16-year-old statistical inference into a confirmed
cosmological systematic. A null result over six sightlines instead
places a $2\sigma$ constraint on the M\'enard prediction, an equally
informative outcome for Stage~IV cosmology. HST/COS is uniquely
required: no other facility delivers the NUV wavelength coverage,
spectral resolution ($R\sim16{,}000$--20{,}000), and UV sensitivity
needed to detect $W_0(2796)\geq 0.3$~\AA\ at
$m_{\rm UV}\sim20.5$--21.5~AB.

\end{document}
